\documentclass{beamer}

\input{preamble.tex}

\title{Instrumental Variables}
\subtitle{Session 2: Interpretation, Formula Instruments, and Applications}
\author{Jingxuan Du \\ \smallskip {\small Dongbei University of Finance and Economics (DUFE)} \\ \smallskip {\small \href{mailto:dujingxuan6@dufe.edu.cn}{dujingxuan6@dufe.edu.cn}}}
\date{Empirical Methods for PhD Students}

\begin{document}

\imageframe{./lecture_includes/cover_frontiers.png}

\section{Introductions}

\begin{frame}{Session 2: Overview}

\textbf{Welcome back!} In Session 1, we covered:
\begin{itemize}
  \item IV mechanics: just-identified IV, 2SLS, weak instruments
  \item Applications: ADH (2013) China Syndrome and BDQ (2020) China roads
\end{itemize}\medskip\pause

\textbf{Session 2 agenda:}
\begin{enumerate}
  \item \textbf{IV Interpretation:} What does IV \emph{really} identify?
  \begin{itemize}
    \item LATE, compliers, treatment effect heterogeneity
    \item Connections to diff-in-diff
  \end{itemize}
  \item \textbf{Formula Instruments:} Shift-share and recentered IV
  \begin{itemize}
    \item Formal conditions, implementation, China HSR illustration
  \end{itemize}
  \item \textbf{Chinese Economics Applications: Session II}
  \begin{itemize}
    \item Meng, Qian, Yared (2015, ReStud): Causes of China's Great Famine
  \end{itemize}
\end{enumerate}

\end{frame}

\begin{frame}{Review: Key Takeaways from Session 1}

Three IV estimands/estimators to know:\medskip
\begin{itemize}
  \item \bgVioletRed{Just-identified IV} $= \rho/\pi$: reduced form over first stage
  \item \bgPictonBlue{2SLS}: optimal combination of multiple instruments; avoids many-IV bias if few instruments used
  \item \bgSun{Wald estimand}: for binary $Z$, $= E[Y|Z=1]-E[Y|Z=0]$ over $E[D|Z=1]-E[D|Z=0]$
\end{itemize}\medskip\pause

Two papers illustrated the ``double-level'' logic for instrument validity:
\begin{itemize}
  \item \textbf{ADH (2013)}: Chinese import shocks in other countries $\rightarrow$ shift-share IV
  \item \textbf{BDQ (2020)}: 1962 historical plan $\rightarrow$ straight-line IV for Chinese roads
\end{itemize}\medskip\pause

A key question we did \emph{not} answer: under what conditions does $\beta^{IV}$ have a \textbf{causal interpretation}? $\rightarrow$ That's our first topic today

\end{frame}

\section{IV Interpretation}

\subsection{LATE and Generalizations}

\begin{frame}{What Does IV Identify, Really?}

IV was invented in the context of structural economic models, typically with a single parameter $\beta$ linearly relating $Y_i$ to $X_i$ 
\begin{itemize}
\item These days we understand $Y_i=\beta X_i+\varepsilon_i$ as describing a \emph{causal} relationship, imposing a (strong) constant-effects restriction
\end{itemize} \pause{}\medskip

The Imbens-Angrist LATE result revolutionized our understanding of IV estimands, and clarified some subtle points around IV identification
\begin{itemize}
\item $\beta^{IV}$ often identifies a convex average of heterogeneous effects under first-stage \emph{monotonicity}: $Z_i$ only affects $X_i$ in one direction
\item IV ``exogeneity'' can arise from two conceptually different assumptions of instrument \emph{independence} and \emph{exclusion}
\end{itemize}
\end{frame}

\begin{frame}{Basic (Binary Treatment, Binary Instrument) Setup}
Let $Y_i(0)$ and $Y_i(1)$ denote individual $i$'s potential outcomes given a binary treatment $D_i\in\{0,1\}$\smallskip
\begin{itemize}
\item Observed outcomes: $Y_i=\left(Y_i(1)-Y_i(0)\right)D_i+Y_i(0)\pause{}\equiv\beta_i D_i+\varepsilon_i$
\end{itemize}\medskip\pause{}
Let $D_i(0)$ and $D_i(1)$ denote individual $i$'s potential treatment given a binary instrument $Z_i\in\{0,1\}$\smallskip
\begin{itemize}
\item Observed treatment: $D_i=\left(D_i(1)-D_i(0)\right)Z_i+D_i(0)$
\end{itemize}\medskip\pause{}
Under what assumptions can we causally interpret \code{ivreg2 Y (D=Z)}?
\end{frame}

\begin{frame}{Imbens and Angrist (1994) Assumptions}
\begin{enumerate}
\item \emph{As-good-as-random assignment}: $Z_i\perp (Y_i(0),Y_i(1),D_i(0),D_i(1))$\smallskip
\begin{itemize}
\item Consider the Angrist draft lottery, or Angrist-Krueger's QoB IV
\end{itemize}\medskip\pause{}
\item \emph{Exclusion}: $Z_i$ only affects $Y_i$ through its effect on $D_i$\smallskip
\begin{itemize}
\item Implicit in the potential outcome notation: $Y_i(d)$ is not indexed by $Z_i$
\end{itemize}\medskip\pause{}
\item \emph{Relevance}: $Z_i$ is correlated with $D_i$\smallskip
\begin{itemize}
\item Equivalently, given Assumption 1, $E[D_{i}(1)-D_i(0)]\neq 0$
\end{itemize}\medskip\pause{}
\item \emph{Monotonicity}: $D_{i}(1)\ge D_{i}(0)$ for all $i$ (i.e., almost-surely)\smallskip
\begin{itemize}
\item The instrument can only shift the treatment in one direction
\end{itemize}
\end{enumerate}
\end{frame}

\begin{frame}{Local Average Treatment Effect (LATE) Identification}
Imbens and Angrist showed that under these assumptions:
\begin{align*}
\beta^{IV}=E[Y_i(1)-Y_i(0)\mid D_i(1)>D_i(0)]
\end{align*}
\vspace{-0.6cm}

The IV estimand $\beta^{IV}$ identifies a LATE: the average treatment effect $Y_i(1)-Y_i(0)$ among \emph{compliers}: those with $1=D_i(1)>D_i(0)=0$\smallskip\pause{}
\begin{itemize}
\item Intuitively, IV can't tell us anything about the treatment effects of \emph{never-takers} $D_i(1)=D_i(0)=0$ or \emph{always-takers} $D_i(1)=D_i(0)=1$\smallskip\pause{}
\item Monotonicity rules out the presence of \emph{defiers}, with $D_i(1)<D_i(0)$
\end{itemize}
\end{frame}

\imageframe{./lecture_includes/angrist_1990.png}

\begin{frame}{What Does This Mean \emph{Practically}?} 
Two conceptually distinct considerations: \emph{internal} vs. \emph{external} validity\smallskip
\begin{itemize}
\item Context of an IV, and who the compliers likely are, may matter\smallskip
\item Usual ``overidentification'' test logic fails: two valid IVs may have different estimands
\end{itemize}\medskip\pause{}
In addition to as-good-as-random assignment / exclusion, we may need to worry about monotonicity when we do IV\smallskip
\begin{itemize}
\item Sensible in many settings (e.g. draft lottery / cancer screening RCTs) \smallskip
\item Maybe questionable in judge IVs (compliers defined by being assigned to a lenient vs. strict judge)
\end{itemize}\medskip\pause{}

Note: technically, we don't need monotonicity if effects are homogenous
\end{frame}

\begin{frame}{Generalizations and Limitations}

\vspace{-0.2cm}
The core logic of IA'94 extends to multivalued treatments/instruments
\begin{itemize}
\item IV identifies an avg. of incremental treatment effects, putting more weight on margins where the instrument shifts the treatment more
\item Regressions of $\mathbf{1}[X_i\ge x]$ on $Z_i$ identify the weights
\end{itemize}\medskip\pause{}
The result does \emph{not} easily extend, however, to multiple treatments 
\begin{itemize}
\item Contamination bias: IV coefficient on $X_{ij}$ generally incorporates effects from other treatments $X_{ik}$, $k\neq j$ 
\end{itemize}\medskip\pause{}
Since 2SLS is a weighted average of ``one-at-a-time'' IVs, overidentified single-treatment specifications can have a LATE interpretation
\begin{itemize}
\item Need all individual IVs to have a LATE interpretation and the 2SLS weights to be convex (the latter can be checked empirically)
\end{itemize}
\end{frame}

\begin{frame}{Generalizations and Limitations (Cont.)}

We can also have a LATE interpretation in IV specifications w/controls
\begin{itemize}
\item E.g. if $Z_i$ is only as-good-as-randomly assigned conditional on $W_i$
\end{itemize}\medskip\pause{}

Key condition: the controls are flexible enough to make $E[Z_i\mid W_i]$ linear
\vspace{-0.7cm}
\begin{itemize}
\item E.g. $W_i$ is a set of mutually exclusive strata dummies
\item IV then identifies a weighted-average of conditional-on-$W_i$ IVs
\end{itemize}\medskip\pause{}

If $E[Z_i\mid W_i]$ is nonlinear, IV may be biased (even w/constant effects)
\begin{itemize}
\item Key practical takeaway: it's good to check sensitivity to how controls are parameterized (add interactions, higher-order polynomials, etc.)
\end{itemize}
\end{frame}

\subsection{Characterizing Compliers}

\begin{frame}{Who Are the Compliers?}
Characterizing the $i$ that make up the IV estimand (w/$D_i(1)>D_i(0)$) is key for understanding internal vs. external validity \smallskip
\begin{itemize}
\item Unfortunately we can't identify compliers directly: we only observe $D_i(1)$ (when $Z_i=1$) or $D_i(0)$ (when $Z_i=0$), not both together!
\end{itemize}\bigskip\pause{}
As it turns out, we can still characterize compliers by their outcomes ($Y_i(0)$ and $Y_i(1)$) and by other observables $X_i$\smallskip
\begin{itemize}
\item Comparing $E[X_i\mid D_i(1)>D_i(0)]$ to $E[X_i]$ can shed light on how $E[Y_i(1)-Y_i(0)\mid D_i(1)>D_i(0)]$ compares to $E[Y_i(1)-Y_i(0)]$
\end{itemize}
\end{frame}

\begin{frame}{Outcomes}
Computing $E[Y_i(1)\mid D_i(1)>D_i(0)]$ is surprisingly easy in the IA setup\smallskip
\begin{itemize}
\item Define $W_i=Y_iD_i$, and note that this new outcome has potentials with respect to $D_i$ of $W_i(1)=Y_i(1)$ and $W_i(0)=0$\smallskip\pause{}
\item Thus IV with $W_i$ as the outcome identifies $E[W_i(1)-W_i(0)\mid D_i(1)>D_i(0)]=E[Y_i(1)\mid D_i(1)>D_i(0)]$
\end{itemize}\medskip\pause{}
Similar logic shows that IV with $Y_i(1-D_i)$ as the outcome and $1-D_i$ as the treatment identifies $E[Y_i(0)\mid D_i(1)>D_i(0)]$\smallskip
\begin{itemize}
\item So easy to do! And extends to covariates / multiple IVs...
\end{itemize}
\end{frame}

\begin{frame}{Illustration: Angrist et al. (2013) Charter School IV}

\begin{center}
	\includegraphics[width=1\textwidth]{./lecture_includes/apw_2013.png}
\end{center}
Decomposing
\vspace{-0.5cm}
\begin{align*}
LATE=&\underbrace{E[Y_i(1)\mid D_i(1)>D_i(0)]-E[Y_i(0)\mid D_i=0]}_{\lambda_1} \\
&- \underbrace{(E[Y_i(0)\mid D_i(1)>D_i(0)]-E[Y_i(0)\mid D_i=0])}_{\lambda_0}
\end{align*}
\vspace{-0.7cm}

shows that charter compliers have typical counterfactual achievement

\end{frame}

\begin{frame}{Covariates}
For covariates $X_i$ (not affected by $D_i$) we can follow a similar trick:
\begin{itemize}
\item Either IV'ing $X_iD_i$ on $D_i$ or IV'ing $X_i(1-D_i)$ on $1-D_i$ identifies complier characteristics $E[X_i\mid D_i(1)>D_i(0)]$\smallskip
\item Shouldn't be very different (implicit balance test); can be averaged
\end{itemize}\bigskip\pause{}
Can be useful to compare with always- and never-taker means:
\begin{itemize}
\item AT: $E[X_i\mid D_i(1)=D_i(0)=1]=E[X_i\mid D_i=1,Z_i=0]$\smallskip
\item NT: $E[X_i\mid D_i(1)=D_i(0)=0]=E[X_i\mid D_i=0,Z_i=1]$
\end{itemize}
\end{frame}

\begin{frame}{Illustration: Angrist et al. (2023) Charter School IV --- Characteristics}

\vspace{-0.2cm}
\begin{center}
	\includegraphics[width=0.9\textwidth]{./lecture_includes/hoee_characteristics.png}
\end{center}

\end{frame}

\begin{frame}{Fancier Things}

General result: $E[g(X_i,Y_i(d))\mid D_i(1)>D_i(0)]$ for any $g(\cdot)$ and any $d\in\{0,1\}$ is identified by $\beta$ in the IV regression:
\begin{align*}
g(X_i,Y_i)\times \mathbf{1}[D_i=d]&=\alpha + \beta\mathbf{1}[D_i=d] + \varepsilon_i\\
\mathbf{1}[D_i=d]&= \mu + \pi Z_i + \nu_i
\end{align*}\pause{}
\vspace{-0.5cm}

Lots of fun stuff we can do here. E.g.: distributions\smallskip
\begin{itemize}
\item $g(X_i,Y_i)=\mathbf{1}[Y_i\le y]$ estimates the CDF of complier potential outcomes, $F(y)=Pr(Y_i(d)<y\mid D_i(1)>D_i(0))$\smallskip
\item $g(X_i,Y_i)=\frac{1}{h}K(\frac{Y_i-y}{h})$ estimates the corresponding PDF, where $K(\cdot)$ is a kernel function and $h$ is a bandwidth 
\end{itemize}

\end{frame}

\begin{frame}{Illustration: Angrist et al. (2023) Charter School IV --- Distributions}
\vspace{-0.2cm}
\begin{center}
	\includegraphics[width=0.9\textwidth]{./lecture_includes/hoee_distribution.png}
\end{center}
\end{frame}

\begin{frame}{Outside the Basic IA Setup}

The same logic applies to IV regressions with flexible controls \smallskip
\begin{itemize}
\item E.g. \emph{ivreg}ing $Y_iD_i$ on $D_i$ instrumenting by $Z_i$ and controlling for cell FE ID's a weighted avg of within-cell $E[Y_i(1)\mid D_i(1)>D_i(0)]$\smallskip\pause{}
\item Abadie (2003) ``kappa-weighting'' is an alternative approach for unweighted averages (and other things)
\end{itemize}\bigskip\pause{}

Logic also goes through with continuous instruments\smallskip
\begin{itemize}
\item E.g. can always mechanically decompose an IV on a binary $D_i$ into implied ``$Y(1)$'' and ``$Y(0)$'' terms
\item Unfortunately, things get trickier with non-binary treatments
\end{itemize}

\end{frame}

\subsection{Diff-in-Diff and IV}

\begin{frame}{What about Diff-in-Diff?}
Most LATE-style discussions of IV identification are ``design-based'': \\ i.e., the instrument is assumed to be drawn randomly, as if in an RCT 
\begin{itemize}
\item But not all instruments are easily thought of this way (e.g. distance)
\end{itemize}\medskip\pause{}

Luckily, whenever you have a reduced form \& first stage that are causal and exclusion+monotonicity hold, the IV ratio has a LATE interpretation
\begin{itemize}
\item E.g. diff-in-diff IV, for distance measure $Z_i\in\{0,1\}$ and $t\in\{1,2\}$:
\vspace{-0.4cm}
\begin{align*}
Y_{it}&=\alpha_i + \tau_t + \beta D_{it} + \epsilon_{it} \\
D_{it}&=\mu_i + \lambda_t + \pi Z_{i}\times\mathbf{1}[t=2] + \epsilon_{it}
\end{align*}
\vspace{-1cm}

Parallel trends in $Z_i$ make the reduced form \& first stage causal\pause{}
\item Recent TWFE literature becomes relevant with fancier specifications
\end{itemize}
\end{frame}

\begin{frame}{Example: Duflo (2001) School-Building Diff-in-Diff}
\begin{center}
	\includegraphics[width=1\textwidth]{./lecture_includes/duflo_2001.png}
\end{center}
Here $Z_i=$ growing up in a region with intensive school building
\end{frame}

\section{Formula Instruments}

\begin{frame}{Exogenous Shocks with Non-Random Exposure}
Increasingly, researchers construct instruments from multiple sources of variation --- only some of which is viewed as exogenous. E.g.:
\begin{itemize}
\item Shift-share instruments $Z_i=\sum_k s_{ik}g_k$, which average together a set of shocks $g_k$ with predetermined exposure shares $s_{ik}$
\item IVs capturing spillovers across social networks/geographies 
\item ``Simulated'' instruments, capturing eligibility for a public program
\end{itemize}\medskip\pause{}

Generically, $Z_i=f_i(g,s)$ where $g=(g_1,\dots,g_k)$ is a set of shocks, $s$ is some measure of shock exposure, and $f_i(\cdot)$ is a known formula
\begin{itemize}
\item How can we leverage exogeneity in $g$, allowing $s$ to be non-random?
\end{itemize}

\end{frame}

\subsection{Shift-Share IV}

\begin{frame}{Example: Autor, Dorn, and Hanson (2013/2014)}
ADH study the effects of rising Chinese import competition on US commuting zones in the 1990's and 2000's
\begin{itemize}
\item Treatment $X_{i}$: growth of Chinese imports in CZ $i$
\item Main outcome $Y_{i}$: change in manufacturing employment in CZ $i$
\end{itemize}\medskip\pause{}

They use a shift-share instrument $Z_{i}=\sum_k s_{ik}g_{k}$, where:
\begin{itemize}
\item Shocks $g_k$: industry $k$'s Chinese import growth in non-US economies
\item Shares $s_{ik}$: pre-period employment share of industry $k$ in CZ $i$
\end{itemize}\medskip\pause{}

Idea: $Z_i$ predicts $X_i$ with proxies for underlying productivity shocks  
\begin{itemize}
\item Fairly intuitive ... but what do we need for it to work? 
\end{itemize}
\end{frame}

\begin{frame}{SSIV with Exogenous Shocks}

Borusyak, Hull, and Jaravel (2022) show shift-share IV can work when
\begin{itemize}
\item The shocks $g_k$ are exogenous, at least conditional on some $q_k$
\item There is enough independent variation in the shocks (i.e. large $K$)
\end{itemize}\bigskip\pause{}

Intuitively, $Z_i=\sum_k s_{ik}g_{k}$ works as a ``translation device,'' bringing the $k$-level variation to bear on causal effects at the $i$ level
\begin{itemize}
\item E.g. in ADH, an industry-level natural experiment can be used to estimate effects at the commuting zone level
\item The shares $s_{ik}$ used for this translation do not need to be exogenous!
\end{itemize}

\end{frame}

\begin{frame}{Two Practical Considerations}

Need to control for $E[Z_i\mid s,q]=E[\sum_k s_{ik}g_{k}\mid s,q]=\sum_k s_{ik}E[g_k\mid q]$ \\ in order to isolate the as-good-as-random variation in shocks
\begin{itemize}
\item When $E[g_k\mid q]=q_k^\prime\gamma$, this means controlling for $\sum_k s_{ik}q_k$ 
\item E.g. in ADH, to isolate within-period shock variation, need to control for $\sum_k s_{ik}\times (\text{period FE})$
\end{itemize}\bigskip\pause{}

Need to adjust SEs for a non-standard clustering: $Z_i=\sum_k s_{ik}g_k$ and $Z_j=\sum_k s_{jk} g_k$ are correlated via their common exposure to shocks
\begin{itemize}
\item Easy fix: use the \emph{ssaggregate} Stata/R packages to translate SSIV estimation to the level of identifying variation (shocks)
\end{itemize}
\end{frame}

\begin{frame}{Estimating Exogenous-Shock SSIV Regressions in Stata}
\vspace{-0.3cm}
\begin{center}
\includegraphics[scale=0.27]{./lecture_includes/ssaggregate.png}
\end{center}
To install: \code{ssc install ssaggregate}
\end{frame}

\begin{frame}{Autor et al. (2013) Revisited}
\begin{center}
\includegraphics[scale=0.52]{./lecture_includes/adh_bhj_restud.png}
\end{center}
\begin{itemize}
\item Columns 3-7 include the key  $\sum_k s_{ik}\times (\text{period FE})$ control
\item Standard errors / first stage F-stat.'s computed via \emph{ssaggregate}
\end{itemize}
\end{frame}

\subsection{Recentered IV}
\begin{frame}{Beyond Shift-Share}

Suppose we wanted to estimate the economic impact of new transportation upgrades (e.g. new high speed rail construction)
\begin{itemize}
\item Economic theory tells us to expect spillovers: easier travel increases productivity across all cities $i$ in a country, to different extents
\item E.g. \emph{market access}: $X_i=\sum_j \tau(g,loc_i,loc_j)^{-1}pop_j$ where $g$ captures the railroad network, $loc_i$ gives lat/lon of cities, and $pop_i$ is city size
\end{itemize}\bigskip\pause{}

Like a more complex shift-share: $Z_i=f_i(g,s)$ for $s=(loc_i,pop_i)_{i=1}^{N}$
\begin{itemize}
\item Suppose we had a natural experiment in high-speed rail construction; can we ``translate'' these $g$ shocks via the market access mapping?
\end{itemize}
\end{frame}

\begin{frame}{Recentered IV}

Borusyak and Hull (2023) show how to leverage exogenous $g$ in $f_i(g,s)$:
\begin{enumerate}
\item Specify \emph{counterfactual} exogenous shocks $g^{(1)},\dots,g^{(C)}$ which may well have occurred (e.g. shuffle the timing of new rail construction)
\item Recompute the instrument: $Z_i^{(c)}=f_i(g^{(c)},s)$, for $c=1,\dots,C$
\item Take the average (``expected instrument''): $\mu_i=\frac{1}{C}\sum_i Z_i^{(c)}$
\end{enumerate}\bigskip\pause{}

Either controlling for $\mu_i$ or using the recentered $\tilde{Z}_i=Z_i-\mu_i$ avoids any bias from endogeneity/non-randomness in $Z_i$
\begin{itemize}
\item Analogous to controlling for $\sum_k s_{ik}q_k$ in shift-share IV! 
\item BH also discuss the analog of non-standard SSIV clustering:\\ the solution is a bit more non-standard (randomization inference)
\end{itemize}

\end{frame}

\begin{frame}[t]{Illustration: High-Speed Rail in China} 
\vspace{-0.3cm}
	\begin{center}
		149 lines were built or planned (as of April 2019)

		\includegraphics[trim={1cm 0.5cm 0.5cm 1cm},clip,width=11cm]{lecture_includes/Lines_actual_planned_nocolor.png}
	\end{center}
\end{frame}

\begin{frame}[t]{Illustration: High-Speed Rail in China} 
\vspace{-0.3cm}
	\begin{center}
		The 83 lines actually built by 2016. Suppose exact timing is random

		\includegraphics[trim={1cm 0.5cm 0.5cm 1cm},clip,width=11cm]{lecture_includes/Line_panel2016.png}
	\end{center}
\end{frame}

\begin{frame}[t]{Illustration: High-Speed Rail in China} 
\vspace{-0.3cm}
	\begin{center}
		A counterfactual draw of 83 lines by 2016

		\includegraphics[trim={1cm 0.5cm 0.5cm 1cm},clip,width=11cm]{lecture_includes/Sim_Line_nlink2016.png}
	\end{center}
\end{frame}

\begin{frame}[t]{Illustration: High-Speed Rail in China} 
\vspace{-0.3cm}
	\begin{center}
		Expected MA growth, $\mu_i$

		\includegraphics[trim={1cm 0.5cm 0.5cm 1cm},clip,width=11cm]{lecture_includes/DateExpected2016.png}
	\end{center}
\end{frame}

\begin{frame}{Illustration: High-Speed Rail in China}
	\begin{center}
Recentered MA growth, $\tilde{Z}_i$

	\includegraphics[trim={1cm 0cm 0cm 1cm},clip,width=11cm]{lecture_includes/NlinkRecentered2016.png}
	\end{center}
\end{frame}

\begin{frame}[label=HSRTable]{Adjusted Estimates of Market Access Effects}
	\begin{center}
	\includegraphics[width=0.9\textwidth]{lecture_includes/hsr_tab.png}
	
	\footnotesize{Regressions of log employment growth on log market access growth in 2007--2016. Spatial-clustered standard errors in parentheses; permutation-based 95\% CI in brackets}
	\end{center}
\end{frame}

\begin{frame}{Recentering for Power}

Leveraging non-random exposure can dramatically improve precision when leveraging exogenous shocks in an IV...
\begin{itemize}
\item ... as long as you recenter to use this variation ``safely''
\end{itemize}\medskip\pause{}

Ex: Medicaid eligibility as a treatment combining statewide policy shocks and individual exposure (income, family structure, etc)
\begin{itemize}
\item In settings where policy shocks are plausibly exogenous, standard approach is to use them directly as instruments (``simulated IV'')
\item BH approach: use eligibility itself, but recenter: e.g. adjust for $i$'s avg. eligibility across permutations of policies (swap MA \& RI, say)
\end{itemize}

\end{frame}

\begin{frame}{Illustration: ACA Medicaid Expansion}

	\begin{center}
	\includegraphics[width=1\textwidth]{lecture_includes/aca_ss.png}
	\end{center}

\footnotesize{1\% ACS sample of non-disabled adults in 2013--14, diff-in-diff IV regressions using one of the two instruments. Controls include state and year fixed effects and an indicator for Republican governor interacted with year. State-clustered standard errors in parentheses; wild score bootstrap 95\% CI in brackets} 

\end{frame}

\section{Chinese Economic Applications: Session II}

\begin{frame}{IV Methods in Chinese Economic History}

Beyond trade and infrastructure, China's unique institutional history offers compelling IV settings:\pause

\begin{itemize}
  \item The Great Leap Forward (1958--62): radical collectivization and grain procurement policies created quasi-exogenous variation in famine severity
  \item One-child policy: regional variation in enforcement creates IVs for fertility
  \item SOE reforms: variation in privatization timing and intensity
  \item Financial repression: capital controls and credit allocation
\end{itemize}\medskip\pause

In this section, we examine a landmark paper using IV to study one of the most consequential episodes in modern Chinese history:

\begin{itemize}
  \item \textbf{Meng, Qian, Yared (2015, ReStud)}: The institutional causes of China's Great Famine, 1959--1961
\end{itemize}

\end{frame}

\begin{frame}{Paper 3: ``The Institutional Causes of China's Great Famine''}

{\small
\textbf{Citation:} Meng, X., Qian, N., and Yared, P. (2015). ``The Institutional Causes of China's Great Famine, 1959--1961.'' \emph{Review of Economic Studies}, 82(4): 1568--1611.
}\medskip\pause

\textbf{Historical context:} The Great Leap Forward Famine (1959--61) is estimated to have killed \textbf{15--55 million people} --- one of the deadliest famines in human history.\medskip\pause

\textbf{Competing explanations:}
\begin{enumerate}
  \item \textbf{Weather} (\emph{natural} cause): Bad harvests due to drought and floods
  \item \textbf{Institutional} (\emph{policy} cause): Grain procurement quotas that extracted food from rural areas even when harvests were poor
\end{enumerate}\medskip\pause

\textbf{Why it matters:} Understanding the cause has implications for policy accountability, state capacity research, and the political economy of authoritarianism

\end{frame}

\begin{frame}{MQY 2015: Research Question and Setting}

\textbf{Central Question:} Were the deaths caused primarily by weather shocks (natural disaster) or by government grain procurement policy (institutional failure)?\medskip\pause

\textbf{Empirical challenge:}
\begin{itemize}
  \item Provinces that had bad harvests \emph{also} received different procurement quotas
  \item Procurement levels and mortality are jointly determined: it's hard to disentangle the causal chain
  \item Selection: provinces with different political environments may respond differently to shocks
\end{itemize}\medskip\pause

\textbf{Key data challenge:} Historical Chinese provincial statistics were suppressed or manipulated during the famine period --- requires careful reconstruction

\end{frame}

\begin{frame}{MQY 2015: IV Strategy}

\textbf{Two-pronged IV approach:}\medskip\pause

\textbf{(1) Weather as IV for grain production:}
\begin{itemize}
  \item Use \textbf{rainfall and temperature} as instruments for actual grain production
  \item Weather is plausibly exogenous to government procurement policy
  \item Identifies the ``weather effect'': how much of mortality variation is explained by food supply shocks alone
\end{itemize}\medskip\pause

\textbf{(2) Grain production shock as IV for procurement:}
\begin{itemize}
  \item The grain production shock (instrumented by weather) shifts procurement quotas
  \item Allows estimation of the effect of \emph{procurement} on mortality, holding production constant
  \item Key identification insight: procurement responded to \emph{planned} output, not actual output --- creating a wedge between food available and food extracted
\end{itemize}

\end{frame}

\begin{frame}{MQY 2015: Data and Methodology}

\textbf{Unit of observation:} 28 Chinese provinces, annual panel 1954--1966\medskip\pause

\textbf{Outcome variables:}
\begin{itemize}
  \item Province-level mortality rates (death rates per 1,000 persons)
  \item Excess mortality relative to pre-famine trend
  \item Grain production, procurement quantities, and grain consumption
\end{itemize}\medskip\pause

\textbf{Data sources:}
\begin{itemize}
  \item Reconstructed provincial mortality statistics from Chinese archival records
  \item Grain production and procurement data from NBS and provincial records
  \item Historical weather station data: rainfall and temperature by province
  \item Political economy variables: provincial party leadership characteristics
\end{itemize}\medskip\pause

\textbf{Specification:} Panel IV regression with province and year fixed effects; two-stage least squares using weather as instrument for production/procurement

\end{frame}

\begin{frame}{MQY 2015: Key Results}

\textbf{Decomposing famine mortality:}\pause
\begin{itemize}
  \item \textbf{Weather} explains only \textbf{12--15\%} of excess mortality variation across provinces
  \item The majority of variation is explained by \textbf{institutional} factors
\end{itemize}\medskip\pause

\textbf{Role of grain procurement:}
\begin{itemize}
  \item A 10\% increase in grain procurement $\rightarrow$ approximately \textbf{4--5\%} increase in the mortality rate
  \item Procurement continued at high levels even as grain production fell sharply
  \item Provinces with higher procurement intensity faced far greater mortality
\end{itemize}\medskip\pause

\textbf{Political economy mechanisms:}
\begin{itemize}
  \item Provinces led by more politically ambitious officials had higher procurement
  \item Information suppression: local officials hid evidence of famine to avoid political censure, exacerbating the crisis
\end{itemize}

\end{frame}

\begin{frame}{MQY 2015: Discussion and Implications}

\textbf{Methodological lessons:}
\begin{itemize}
  \item Illustrates IV in a \textbf{historical/archival} setting with imperfect data
  \item The weather instrument satisfies: relevance (weather $\rightarrow$ grain production), exclusion (weather affects mortality only through food supply)
  \item Province fixed effects absorb time-invariant confounders; year FEs absorb common shocks
\end{itemize}\medskip\pause

\textbf{Broader implications:}
\begin{itemize}
  \item Famines in democratic societies are rare (Sen 1981): accountability mechanisms prevent government inaction
  \item Authoritarian information suppression can turn a natural shock into a catastrophe
  \item Evidence that institutional design, not just natural endowments, drives development outcomes
\end{itemize}\medskip\pause

\textbf{Legacy:} This paper spawned a large literature on the political economy of China's Great Leap Forward and on information and governance in autocracies

\end{frame}

\begin{frame}{Summary: Chinese Economics Papers Using IV}

\begin{center}
\small
\begin{tabular}{p{1.5cm}p{2.3cm}p{2.0cm}p{2.5cm}}
\hline
\textbf{Paper} & \textbf{Question} & \textbf{IV} & \textbf{Finding} \\
\hline
ADH (2013) \newline \emph{AER} & China trade $\rightarrow$ US jobs & Imports in other countries (shift-share) & $-$2M mfg. jobs 1990--2007 \\[4pt]
BDQ (2020) \newline \emph{JPE} & Roads $\rightarrow$ China GDP & Historical plan (straight-line) & +10--15\% growth near highways \\[4pt]
MQY (2015) \newline \emph{ReStud} & Causes of Great Famine & Weather $\rightarrow$ production $\rightarrow$ procurement & Policy drove 85\%+ of deaths \\
\hline
\end{tabular}
\end{center}\medskip\pause

\textbf{Common themes:}
\begin{itemize}
  \item All three papers use creative IV designs to solve sharp endogeneity problems
  \item Exogenous variation comes from unexpected sources: competitor imports, military plans, weather
  \item Results have major policy implications for trade, development, and governance
\end{itemize}

\end{frame}

\section{Concluding Thoughts}

\begin{frame}{Concluding Thoughts}

We've covered a substantial amount of ground in these two sessions!\pause
\begin{itemize}
  \item From basic IV mechanics to advanced recentering and formula instruments
  \item From classic Angrist (1990) design to cutting-edge Chinese economic applications
\end{itemize}\bigskip\pause{}

Some general advice on how to approach IV in the ``real world'':
\begin{enumerate}
\item Figure out what the reduced form and first stage regressions are, \\ \& what variation in the instrument(s) they leverage via the controls
\item Assess the plausibility of as-good-as-random assignment (or parallel trends) of this variation, both theoretically and empirically
\item Interrogate exclusion, which lets you interpret RF/FS, and monotonicity/complier stats if heterogeneous effects seem important
\end{enumerate}
\end{frame}

\begin{frame}{General Advice for PhD Students}

\textbf{Finding a good IV:}\pause
\begin{itemize}
  \item Think about your setting's institutional details --- what creates exogenous variation?
  \item Classic sources: lotteries, policy discontinuities, historical accidents, geographic features
  \item For China specifically: policy rollout timing, administrative boundaries, historical shocks
\end{itemize}\medskip\pause

\textbf{Evaluating IV validity:}
\begin{itemize}
  \item Assess relevance with first-stage F-statistics ($>10$ rule of thumb)
  \item Check exclusion with falsification tests, balance tests, and robustness checks
  \item Think carefully about who the compliers are and whether LATE is the object of interest
\end{itemize}\medskip\pause

\textbf{Practical tools:}
\begin{itemize}
  \item Stata: \code{ivreg2}, \code{ssaggregate}, \code{weakivtest}
  \item R: \code{AER::ivreg()}, \code{fixest::feols()}, \code{ssaggregate}
\end{itemize}

\end{frame}

\begin{frame}{Further Resources}

\textbf{Textbooks and reviews:}
\begin{itemize}
  \item Angrist and Pischke (2009): \emph{Mostly Harmless Econometrics}
  \item Angrist and Pischke (2015): \emph{Mastering 'Metrics}
  \item Imbens and Rubin (2015): \emph{Causal Inference in Statistics, Social, and Biomedical Sciences}
\end{itemize}\medskip\pause

\textbf{Methodological frontiers:}
\begin{itemize}
  \item Borusyak, Hull, Jaravel (2022): Shift-share IV formal conditions
  \item Borusyak and Hull (2023): Recentered IV
  \item Goldsmith-Pinkham, Sorkin, Swift (2020): SSIV shares as IVs
  \item Angrist and Koles\'ar (2022): Weak instrument inference
\end{itemize}\medskip\pause

\textbf{China-specific:}
\begin{itemize}
  \item Brandt and Rawski (eds., 2008): \emph{China's Great Economic Transformation}
  \item Lu and Yu (2015): Trade liberalization and markup dispersion in China
\end{itemize}

\end{frame}

\begin{frame}{Your IV Research: Building a Pipeline}

Three questions to ask yourself at every step:\medskip
\begin{enumerate}
  \item \textbf{What is my first stage?} \\
  Is the instrument strongly correlated with the treatment? Check the F-stat.\medskip\pause
  \item \textbf{What is my reduced form?} \\
  Does the instrument predict the outcome? Is the sign sensible?\medskip\pause
  \item \textbf{What is my exclusion restriction?} \\
  Is there any channel by which the instrument could affect the outcome other than through the treatment? Can you test it?
\end{enumerate}\medskip\pause

If you can answer all three convincingly, you likely have a credible IV design --- the foundation of publishable empirical work

\end{frame}

\begin{frame}{Keep Calm and \code{ivreg2} On!}

\begin{center}
Thank you for your attention throughout both sessions!

\bigskip
\textbf{Jingxuan Du} \\
Dongbei University of Finance and Economics (DUFE) \\
\smallskip
\url{dujingxuan6@dufe.edu.cn}

\bigskip
{\small Good luck on your future adventures with IV!}
\end{center}
\end{frame}

\end{document}
